%%%%%%%%%%%%%%%%%%%%%%%%%%%%%%%%%%%%%%%%%%%%%%%%%%%%%%%%%%%%%%%%%%%%%%
% Lớp 10 | Phương pháp tọa độ trong mặt phẳng (𝑂𝑥𝑦) | Tích vô hướng (theo tọa độ)
% Dạng 0H9?2-2: Tìm góc nhờ tích vô hướng
% Số câu: 10  (sinh tự động bởi tach_cau.py; từ nay thầy sửa trực tiếp file này)
%%%%%%%%%%%%%%%%%%%%%%%%%%%%%%%%%%%%%%%%%%%%%%%%%%%%%%%%%%%%%%%%%%%%%%

%%%%%%%%%%%%%%%% TRẮC NGHIỆM 4 PHƯƠNG ÁN %%%%%%%%%%%%%%%%

% nguồn: MaHoaNH_TracNghiem_Lop10
\begin{ex}%[0H9H2-2]%[ma:39b78301]
	Trên mặt phẳng tọa độ $Oxy$, góc giữa hai véc-tơ  $\overrightarrow{a}=(1;-1)$ và $\overrightarrow{b}=(-2;0)$ có số đo bằng
	 \choice
		{$60^{\circ}$} 
		 {$45^{\circ}$}
		{$90^{\circ}$} 
		{\True $135^{\circ}$} 
	\loigiai{ 
		Ta có  $\cos\left(\overrightarrow{a},\overrightarrow{b}\right)=\dfrac{1\cdot(-2)-1\cdot0}{\sqrt{1^2+(-1)^2}\cdot\sqrt{(-2)^2+0^2}}=\dfrac{-1}{\sqrt{2}}$.\\
		Suy ra $\left(\overrightarrow{a},\overrightarrow{b}\right)=135^\circ$. 
	} 
\end{ex}

% nguồn: MaHoaNH_TracNghiem_Lop10
\begin{ex}%[0H9B2-2]%[ma:82126716]
	Trong hệ trục $Oxy$, cho điểm $E(-2;1)$ và véc-tơ $\vec{n}=(4;3)$. Phương trình tổng quát của đường thẳng $d$ đi qua điểm $E$ và  có véc-tơ pháp tuyến $\vec{n}$ là
	\choice 
		 {$-2x+y-5=0$}
		{$4x+3y-5=0$} 
		 {$3x-4y+10=0$}
		{\True $4x+3y+5=0$} 
	\loigiai{ 
		Đường thẳng $d$ nhận $\vec{n}=(4;3)$ làm véc-tơ pháp tuyến nên phương trình đường thẳng  $d$ có dạng $4x+3y+c=0$.\\
		Vì đường thẳng $d$ đi qua điểm $E(-2;1)$, nên ta có:  $4\cdot(-2) +3\cdot1+c=0 \Leftrightarrow c=5$.\\
		Vậy phương trình tổng quát của đường thẳng  $d\colon 4x+3y+5=0$.
	} 
\end{ex}


%%%%%%%%%%%%%%%% ĐÚNG/SAI %%%%%%%%%%%%%%%%

% nguồn: MaHoaNH_DungSai_Lop10
\begin{ex}%[0H9N2-2]%[ma:a7aa0c0e]
	Trong hệ tọa độ $Oxy$, cho các véc-tơ $\overrightarrow{a}=(1;3)$, $\overrightarrow{b}=(-4;3)$. Gọi $\alpha$ là góc giữa hai véc-tơ $\overrightarrow{a}$, $\overrightarrow{b}$. Các  khẳng định sau đây đúng hay sai?
	 \choiceTF
		{\True $\overrightarrow{a}\cdot\overrightarrow{b}=5$} 
		{$\alpha\in  (90^\circ; 180^\circ)$}
		 {\True $\cos\alpha>0$}
		 {\True $\alpha<90^\circ$}
	\loigiai{ 
		Ta có  $\left|\overrightarrow{a}\right|=\sqrt{1^2+3^2}=\sqrt{10}$, $\left|\overrightarrow{b}\right|=\sqrt{(-4)^2+3^2}=5$.\\
		Suy ra $\overrightarrow{a}\cdot\overrightarrow{b}=1\cdot (-4)+3\cdot  3=5$.\\
		 Do đó $\cos\alpha=\dfrac{\overrightarrow{a}\cdot\overrightarrow{b}}{\left|\overrightarrow{a}\right|\cdot \left|\overrightarrow{b}\right|}=\dfrac{1}{\sqrt{10}}$.
		Vậy $\alpha\approx  71^\circ$.
	} 
\end{ex}

% nguồn: MaHoaNH_DungSai_Lop10
\begin{ex}%[0H9H2-2]%[ma:0d144d01]
	Trong hệ tọa độ $Oxy$, cho các véc-tơ $\overrightarrow{a}=(4;3)$ và $\overrightarrow{b}=(1;7)$. Trong  các khẳng định sau, khẳng định nào đúng, khẳng định nào sai?
	 \choiceTF
		{\True  $(\overrightarrow{i}; \overrightarrow{j})=90^\circ$}
		{\True  $(\overrightarrow{a}-\overrightarrow{b})\perp \overrightarrow{a}$}
		{$\cos (\overrightarrow{a}; \overrightarrow{i})=\dfrac{3}{5}$} 
		{\True $(\overrightarrow{a};  \overrightarrow{b})=45^\circ$}
	\loigiai{ 
		\begin{itemchoice}
			\itemch Do $\overrightarrow{i}$, $\overrightarrow{j}$ là các véc-tơ đơn vị lần lượt nằm  trên các trục $Ox$, $Oy$ nên $(\overrightarrow{i}; \overrightarrow{j})=90^\circ$.
			 \itemch Ta có $\overrightarrow{a}-\overrightarrow{b}=(3;-4)$.\\
			Suy  ra $(\overrightarrow{a}-\overrightarrow{b})\cdot \overrightarrow{a}=3\cdot 4-4\cdot 3=0$.\\
			Vậy $(\overrightarrow{a}-\overrightarrow{b})\perp \overrightarrow{a}$. 
			\itemch Ta có $\cos  \left(\overrightarrow{a}, \overrightarrow{i}\right)=\dfrac{\overrightarrow{a}\cdot\overrightarrow{i}}{\left|\overrightarrow{a}\right|\cdot\left|\overrightarrow{i}\right|}=\dfrac{4\cdot 1+3\cdot 0}{\sqrt{4^2+3^2}\cdot 1}=\dfrac{4}{5}$.
			\itemch Ta có $\cos \left(\overrightarrow{a}, \overrightarrow{b}\right)=\dfrac{\overrightarrow{a}\cdot\overrightarrow{b}}{\left|\overrightarrow{a}\right|\cdot\left|\overrightarrow{b}\right|}=\dfrac{4\cdot 1+3\cdot  7}{\sqrt{4^2+3^2}\cdot\sqrt{1^2+7^2}}=\dfrac{\sqrt{2}}{2}$.\\
			Suy ra $\left(\overrightarrow{a},\overrightarrow{b}\right)=45^\circ$. 
		\end{itemchoice}
	} 
\end{ex}


%%%%%%%%%%%%%%%% TỰ LUẬN %%%%%%%%%%%%%%%%

% nguồn: MaHoaNH_TuLuan_Lop10
\begin{bt}%[0H9Y2-2]%[ma:e8f3abf9]
	Trong mặt phẳng tọa độ $Oxy$, cho tam giác $ABC$ có $A(-2;1)$, $B(3;0)$, $C(5;-6)$. Viết phương trình đường thẳng  chứa đường trung tuyến $AM$ ($M$ là trung điểm của $BC$).
	 \loigiai{
		Ta có $M$  là trung điểm của $BC$ nên $M(4;-3)$.\\
		Ta cũng có  $\overrightarrow{AM}= ( x_M-x_A;y_M-y_A)=(6;-4)$.\\
		Đường thẳng $AM$ đi qua $A(-2;1)$ có VTCP $\overrightarrow{u}=(6;-4)$ nên có  phương trình là
		\[\heva{&x=-2+6t \\  &y=1-4t} (t\in \mathbb{R})\]
	 }
\end{bt}

% nguồn: MaHoaNH_TuLuan_Lop10
\begin{bt}%[0H9B2-2]%[ma:fa6e5860]
	Trong mặt phẳng $(Oxy)$  cho điểm $M(2;4)$ và đường thẳng $d\colon 3x+y+2022=0$. Viết phương trình đường thẳng $d'$ song song với đường thẳng $d$ và cách điểm $M$ một khoảng bằng $\sqrt{10}$.
	\loigiai{ 
		 Do $d'\parallel d$ nên $d'$ có phương trình dạng $3x+y+c=0$ với $c\ne 2022$.\\
		 Ta có $\mathrm{d}(M,d')=\sqrt{10}$ nên $\dfrac{\left| 6+4+c\right|}{\sqrt{10}}=\sqrt{10}\Leftrightarrow\left| 10+c\right|=10$.\\
		Suy ra $10+c=10$ hoặc  $10+c=-10$ hay $c=0$ hay $c=-20$.\\
		Vậy có  $2$ đường thẳng thoã mãn ycbt là $d'\colon 3x+y=0$; $d'\colon 3x+y-20=0$.
	 }
\end{bt}

% nguồn: MaHoaNH_TuLuan_Lop10
\begin{bt}%[0H9B2-2]%[ma:aa8d60ff]
	Trong mặt phẳng tọa độ $Oxy$, cho điểm $N(-2;-3)$ và đường thẳng $d\colon-3x+5y+9=0$. Viết phương trình của đường thẳng $\Delta$ đi qua $N$ và $\Delta$ vuông góc với đường thẳng $d$. 
	 \loigiai{
		\begin{itemize}
			\item $d$ có véc-tơ  pháp tuyến $\overrightarrow{n}=(-3;5)$.
			\item Vì $\Delta\perp d$ nên $\Delta$ có véc-tơ pháp tuyến  $\overrightarrow{m}=(5;3)$.
			\item Phương trình của đường  thẳng $\Delta$ là
			 $$5(x+2)+3(y+3)=0\Leftrightarrow5x+3y+19=0.$$
		\end{itemize}
	 }
\end{bt}

% nguồn: MaHoaNH_TuLuan_Lop10
\begin{bt}%[0H9B2-2]%[ma:96bacc95]
	Cho ba điểm $A(-3;2)$, $B(3;4)$, $C(9;-2)$. Viết phương trình  tổng quát và phương trình tham số của đường cao $AH$.
	 \loigiai{
		Ta có $\overrightarrow{BC}=(6;-6)$, suy ra véc-tơ cùng phương với $\overrightarrow{BC}$ là $\vec{n}=(1;-1)$.\\ 
		Do $AH\perp BC\Rightarrow\vec{n}=(1;-1)$ là véc-tơ  pháp tuyến của đường cao $AH$.\\
		Phương trình tổng quát của đường cao $AH$ là 
		$$1(x+3)-1(y-2)=0\Leftrightarrow  x-y+5=0.$$
		Vì $\vec{n}=(1;-1)$ là véc-tơ pháp tuyến của đường  cao $AH$, suy ra $\vec{a}=(1;1)$ là véc-tơ chỉ phương của đường cao $AH$.\\
		Phương trình  tham số của đường cao $AH$ là
		$$\heva{&x=-3+t\\&y=2+t}, t\in\mathbb{R}.$$ 
	 }
\end{bt}

% nguồn: MaHoaNH_TuLuan_Lop10
\begin{bt}%[0H9B2-2]%[0H9B3-2]%[ma:281fb5f1]
	Cho $A(2;-1)$; $B(4;  3)$; $C(-5; 0)$
	\begin{listEX}
		\item Viết phương trình tổng quát  đường thẳng đi qua $2$ điểm $A$ và $C$.
		 \item Viết phương trình đường tròn tâm $B$ và tiếp xúc với đường thẳng $AC$.
	\end{listEX}
	 \loigiai{
		\begin{listEX}
			\item Đường thẳng $AC$ có véc-tơ chỉ phương $\overrightarrow{A C}=(-7;  1)$.\\
			Suy  ra đường thẳng $AC$ qua $A$, có véc-tơ pháp tuyến $\vec{n}=(1; 7)$ có phương trình
			\[1\cdot (x-2)+7\cdot (y+1)=0 \Leftrightarrow x+7y+5=0  \]
			\item Ta có đường tròn tâm $B$ và tiếp xúc với đường thẳng  $AC$ nên
			\[R=\mathrm{d}\,(B,AC)=\dfrac{|4+7 \cdot 3+5|}{\sqrt{1^2+7^2}}=\dfrac{30}{\sqrt{50}}=3 \sqrt{2}.\] 
			Suy ra đường tròn $(C)$ có tâm $B$,  bán kính $R=3 \sqrt{2}$ có phương trình
			 \[(C)\colon (x-4)^2+(y-3)^2=18.\]
		\end{listEX}
	} 
\end{bt}

% nguồn: MaHoaNH_TuLuan_Lop10
\begin{bt}%[0H9K2-2]%[ma:df3fea04]
	Trong mặt phẳng  tọa độ $Oxy$, cho $A(-2;1)$, $B(2;-3)$ và đường thẳng $(d)\colon 3x + 4y + 7 =0$.
	\begin{enumerate}
		\item Viết  phương trình tham số của đường thẳng $(d_1)$ đi qua hai điểm $A$, $B$.
		\item Viết phương trình đường thẳng $(d_2)$  đi qua $A$ và vuông góc với đường thẳng $(d)$.
		\item Tính $\cos $ của góc giữa hai đường thẳng  $(d_1)$ và $(d_2)$.
	\end{enumerate}
	 \loigiai{
		\begin{enumerate}
			\item Ta có $\overrightarrow{AB}=(4;-4)$.\\ 
			 Ta có $(d_1)$ qua $A(-2;1)$ và có véc-tơ chỉ phương $\overrightarrow{AB}=(4;-4)$ nên có phương trình tham số là
			\[\heva{&x = -2 + 4t\\&y=1  -4t} \quad (t\in \mathbb{R}).\]
			\item Ta có $(d_2)\perp (d)\colon 3x +  4y + 7=0$ nên $(d_2)\colon 4x - 3y+m=0$.\\
			Mà  $A(-2;1)\in (d_2)$ nên
			\[4\cdot (-2) -3\cdot 1  + m =0 \Leftrightarrow -11 + m =0 \Leftrightarrow m=11.\]
			 Vậy $(d_2)\colon 4x - 3y +11=0$.
			 \item Ta có $(d_1)$ có véc-tơ pháp tuyến là $\overrightarrow{n}_1 = (4;4)$, $(d_2)$ có véc-tơ pháp tuyến là $\overrightarrow{n}_2 = (4;-3)$.\\
			Khi đó 
			\[\cos (d_1,d_2)=\dfrac{|\overrightarrow{n}_1 \cdot \overrightarrow{n}_2|}{\left|\overrightarrow{n}_1\right|\cdot  \left|\overrightarrow{n}_2\right|}=\dfrac{|4\cdot 4 + 4\cdot (-3)|}{\sqrt{4^2 + 4^2}\cdot \sqrt{4^2 + (-3)^2}}=\dfrac{\sqrt{2}}{10}.\]
		\end{enumerate}
	 }
\end{bt}
